\chapter{Цепи Маркова}

Рассмотрим последовательность случайных величин $ \xi_0, \xi_1, \dotsc $, принимающих
значения из некоторого конечного множества $ \Set{x_1, \dots, x_N} $.
Будем для простоты считать, что эти значения совпадают с набором чисел $ \Set{1, \dots, N} $ и
называть их \emph{состояниями}.

\begin{definition}
    Последовательность $ \xi_0, \xi_1, \dotsc $ называется \emph{цепью Маркова с конечным числом
    ($ N $) состояний}, если
    $$ \forall n \quad \forall a_0, a_1, \dotsc \quad
    P\Set{\xi_{n + 1} = a_{n + 1} | \xi_n = a_n, \quad \dots, \quad \xi_0 = a_0} =
    P\Set{\xi_{n + 1} = a_{n + 1} | \xi_n = a_n} $$
    (``будущее при фиксированном настоящем не зависит от прошлого'').
\end{definition}

В определении важную роль играют \emph{переходные состояния} $ P\Set{\xi_{n + 1} = j | \xi_n = i} $.
Если эти вероятности не зависят от момента перехода $ n $, то такая цепь называется
\emph{однородной}.
Каждой однородной цепи Маркова соответствует \emph{матрица переходов за один шаг}
$$ P =
\begin{pmatrix}
    p_{11} & \dots & p_{1N} \\
    \vdots & & \vdots \\
    p_{N1} & \dots & p_{NN}
\end{pmatrix} $$
Мы будем рассматривать только однородные цепи.

Ещё один важный объект "--- \emph{матрица переходов за $ n $ шагов}:
$$ \nder P =
\begin{pmatrix}
    \nder p_{11} & \dots & \nder p_{1N} \\
    \vdots & & \vdots \\
    \nder p_{N1} & \dots & \nder p_{NN}
\end{pmatrix} $$

\begin{remark}
    Любая матрица $ \nder P $ является \emph{стохастической}, то есть
    \begin{enumerate}
        \item $ \nder p_{ij} \ge 0 $;
        \item $ \sum_j \nder p_{ij} = 1 $.
    \end{enumerate}
\end{remark}

\section{Классификация состояний цепей Маркова}

\begin{enumerate}
    \item Состояние $ i $ достижимо из $ j $, если существует такое $ n $, что $ \nder p_{ij} > 0 $.
    \item Состояния $ i $ и $ j $ \emph{сообщаются}, если $ i $ достижимо из $ j $ и $ j $
        достижимо из $ i $.
    \item Состояние $ i $ "--- \emph{существенное}, если оно достижимо из любого состояния,
        достижимого из $ i $.
        То есть, выйдя из этого состояния, всегда можно с ненулевой вероятностью в него вернуться.
    \item Состояние $ i $ "--- \emph{несущественное}, если оно не достижимо из некоторого $ j $,
        достижимого из $ i $.
\end{enumerate}

\begin{remark}
    Все существенные состояния образуют классы эквивалентности.
\end{remark}

\begin{definition}
    Если у ЦМ только один класс существенных состояний, то она называется \emph{неразложимой}.
\end{definition}

\begin{theorem}[уравнение Чепмена"--~Колмогорова]
    $ \nder[k + m]P = \nder[k]P \nder[m]P = \nder[m]P \nder[k]P $
\end{theorem}

\begin{implication}
    $ \nder[n]P = P^n $
\end{implication}

\begin{implication}
    Пусть $ p_n $ "--- вектор, состоящий из вероятностей попасть в какое-то состояние не на первом,
    а на $ n $-м шаге

    $$ p_n = p_0P^n, \quad p_n = p_{n - 1}P $$
\end{implication}

\section{Простейшая эргодическая теорема}

Рассмотрим простейший случай, когда имеем только два состояния цепи, и в матрице переходов за один
шаг все вероятности ненулевые.
Поставим такую задачу: имея матрицу
$$
\begin{pmatrix}
    p_{11} & p_{12} \\
    p_{21} & p_{22}
\end{pmatrix} $$
необходимо найти пределы вероятностей $ \lim P_{i1}(n) $ и $ \lim P_{i2}(n) $, где $ P_{i1}(n) =
P\Set{\xi_n = 1 | \xi_0 = i} $.

Очевидно, что в нашем случае достаточно найти один из пределов, поскольку сумма этих двух пределов
равна 1.
$$ P_{i1}(n + 1) = P_{i1}(n)p_{11} + P_{i2}(n) p_{21} = P_{i1}(n) p_{11} +
\bigl( 1 - P_{i1}(n) \bigr)p_{21} = p_{21} + p_{i1}(n)(p_{11} - p_{21}) $$
$$ \forall c \quad P_{i1}(n + 1) - c = (p_{i1}(n) - c)(p_{11} - p_{21}) + p_{21} - c(p_{12} + p_{21}) $$
Возьмём $ c = \dfrac{p_{21}}{p_{12} + p_{21}} $.
Обозначим $ q = p_{11} - p_{21} $.
$$ P_{i1}(n + 1) - c = q \bigl( P_{i1}(n) - c \bigr) $$
Применим это равенство $ n $ раз.
$$ P_{i1}(n + 1) - c = q^n(p_{i1} - c) \underimp{0 < q < 1} \lim P_{i1}(n) =
\frac{p_{21}}{p_{12} + p_{21}}, \quad \lim P_{12} = \frac{p_{12}}{p_{12} + p_{21}} $$

\section{Возвратные состояния цепи Маркова}

Введём обозначения для вероятностей первого возвращения в какой-то момент в исходное состояние $ j $:
$$ f_j(n) = P\Set{\xi_n = j, \quad \xi_{n - 1} \ne j, \quad \dots, \quad \xi_1 \ne j, \quad
\xi_0 = j}, \quad f_j(0) = 0 $$

\begin{definition}
    Состояние $ j $ называется \emph{возвратным}, если
    $$ \sum_{n = 1}^\infty f_j(n) = 1 $$
    иначе "--- \emph{невозвратным}.
\end{definition}

Введём ещё одно обозначение: $ p_j(n) = P\Set{\xi_n = j | \xi_0 = j} $.

\begin{theorem}[критерий возвратности]
    Состояние $ j $ является возвратным \textbf{тогда и только тогда}, когда
    $$ \sum_{n = 1}^\infty p_j(n) = \infty $$
\end{theorem}

\begin{proof}
    Рассмотрим производящие функции:
    $$ F(s) = \sum_{n = 0}^\infty f_j(n)s^n, \quad P(s) = \sum_{n = 0}^\infty p_j(n) s^n $$

    Заметим, что
    $$ P(s) = \frac1{1 - F(s)} $$
    По свойствам производящих функций
    $$ \sum_{n = 0}^\infty f_j(n) = \lim\limits_{s \to 1} F(s) $$
    $$ \sum_{n = 0}^\infty p_j(n) = \lim\limits_{s \to 1} P(s) $$
\end{proof}
